% ============================================================================
% BỔ SUNG ĐỢT 2 — TOÁN 9 KNTT CHƯƠNG V: BÀI 14 & BÀI 15 (36 CÂU)
% Mỗi mục nhỏ thêm 2 TN Nhận biết + 2 Đúng/Sai Thông hiểu + 2 TLN Vận dụng
% ============================================================================

% --- 9H5?14-1: Quan hệ giữa đường kính và dây ---

\begin{ex}%[9H5N14-1]
% nguon: SGK Toán 9 KNTT tập 1 – Bài 14, Cung và dây của một đường tròn
	Cho đường tròn $(O; 7{,}5\text{ cm})$. Độ dài đường kính của đường tròn đó bằng
	\choice
	{$7{,}5\text{ cm}$}
	{\True $15\text{ cm}$}
	{$22{,}5\text{ cm}$}
	{$30\text{ cm}$}
	\loigiai{
		Độ dài đường kính của đường tròn bằng hai lần bán kính: $d = 2R = 2 \cdot 7{,}5 = 15\text{ cm}$.
	}
\end{ex}

\begin{ex}%[9H5N14-1]
% nguon: SGK Toán 9 KNTT tập 1 – Bài 14, Cung và dây của một đường tròn
	Với mọi dây $AB$ của đường tròn $(O; R)$, khẳng định nào sau đây luôn đúng?
	\choice
	{$AB < R$}
	{$AB = 2R$}
	{\True $AB \le 2R$}
	{$AB \ge 2R$}
	\loigiai{
		Trong một đường tròn, đường kính là dây lớn nhất, do đó mọi dây $AB$ đều thỏa mãn $AB \le 2R$.
	}
\end{ex}

\begin{ex}%[9H5H14-1]
% nguon: SGK Toán 9 KNTT tập 1 – Bài 14, Cung và dây của một đường tròn
	Cho đường tròn $(O; 10\text{ cm})$ có dây $MN = 16\text{ cm}$ và đường kính $PQ$ vuông góc với $MN$ tại điểm $I$ ($P$ gần $I$ hơn $Q$). Xét tính đúng, sai của các khẳng định sau:
	\choiceTF
	{\True Độ dài đường kính $PQ$ bằng $20\text{ cm}$}
	{\True Điểm $I$ là trung điểm của $MN$ và $IM = 8\text{ cm}$}
	{Độ dài đoạn thẳng $OI$ bằng $5\text{ cm}$}
	{\True Diện tích tam giác $PMN$ bằng $32\text{ cm}^2$}
	\loigiai{
		\begin{itemchoice}
			\itemch Đúng. Đường kính $PQ = 2R = 2 \cdot 10 = 20\text{ cm}$.
			\itemch Đúng. Đường kính vuông góc với một dây thì đi qua trung điểm của dây ấy nên $IM = \dfrac{MN}{2} = 8\text{ cm}$.
			\itemch Sai. Trong tam giác $OIM$ vuông tại $I$: $OI = \sqrt{OM^2 - IM^2} = \sqrt{10^2 - 8^2} = 6\text{ cm}$.
			\itemch Đúng. Vì $P$ gần $I$ hơn nên $PI = OP - OI = 10 - 6 = 4\text{ cm}$. Diện tích tam giác $PMN$ là $\dfrac{1}{2} \cdot MN \cdot PI = \dfrac{1}{2} \cdot 16 \cdot 4 = 32\text{ cm}^2$.
		\end{itemchoice}
	}
\end{ex}

\begin{ex}%[9H5H14-1]
% nguon: SGK Toán 9 KNTT tập 1 – Bài 14, Cung và dây của một đường tròn
	Cho tam giác $ABC$ vuông tại $A$ có $AB = 9\text{ cm}$, $AC = 12\text{ cm}$. Gọi $(O)$ là đường tròn đường kính $BC$. Xét tính đúng, sai của các khẳng định sau:
	\choiceTF
	{\True Độ dài cạnh $BC$ bằng $15\text{ cm}$}
	{\True Điểm $A$ nằm trên đường tròn $(O)$}
	{Bán kính đường tròn $(O)$ bằng $8\text{ cm}$}
	{\True Khoảng cách từ tâm $O$ đến dây $AC$ bằng $4{,}5\text{ cm}$}
	\loigiai{
		\begin{itemchoice}
			\itemch Đúng. Áp dụng định lí Pythagore cho tam giác $ABC$ vuông tại $A$: $BC = \sqrt{9^2 + 12^2} = 15\text{ cm}$.
			\itemch Đúng. Vì $O$ là trung điểm cạnh huyền $BC$ nên $OA = OB = OC = \dfrac{BC}{2}$, suy ra $A$ thuộc đường tròn đường kính $BC$.
			\itemch Sai. Bán kính đường tròn $(O)$ là $R = \dfrac{BC}{2} = 7{,}5\text{ cm}$.
			\itemch Đúng. Gọi $M$ là trung điểm $AC$, ta có $OM \perp AC$ và $OM$ là đường trung bình của tam giác $ABC$ nên $OM = \dfrac{AB}{2} = 4{,}5\text{ cm}$.
		\end{itemchoice}
	}
\end{ex}

\begin{bt}%[9H5V14-1]
% nguon: SGK Toán 9 KNTT tập 1 – Bài 14, Cung và dây của một đường tròn
	Cho đường tròn $(O; 15\text{ cm})$. Dây $AB$ vuông góc với đường kính $CD$ tại điểm $H$ sao cho $CH = 3\text{ cm}$. Tính độ dài dây $AB$ (đơn vị cm).
	\shortans{18}
	\loigiai{
		Vì $H$ nằm trên bán kính $OC$ nên $OH = OC - CH = 15 - 3 = 12\text{ cm}$.\\
		Do $CD \perp AB$ tại $H$ nên $H$ là trung điểm của $AB$. Trong tam giác $OHA$ vuông tại $H$:
		$$AH = \sqrt{OA^2 - OH^2} = \sqrt{15^2 - 12^2} = 9\text{ cm}.$$
		Vậy $AB = 2AH = 18\text{ cm}$.
	}
\end{bt}

\begin{bt}%[9H5V14-1]
% nguon: SGK Toán 9 KNTT tập 1 – Bài 14, Cung và dây của một đường tròn
	Cho đường tròn $(O)$ có đường kính $AB = 26\text{ cm}$. Dây $CD$ vuông góc với đường kính $AB$ tại điểm $K$ nằm giữa $O$ và $A$ sao cho $CD = 24\text{ cm}$. Tính độ dài đoạn thẳng $AK$ (đơn vị cm).
	\shortans{8}
	\loigiai{
		Bán kính đường tròn là $R = \dfrac{AB}{2} = 13\text{ cm}$.\\
		Vì $AB \perp CD$ tại $K$ nên $CK = \dfrac{CD}{2} = 12\text{ cm}$.\\
		Trong tam giác $OKC$ vuông tại $K$: $OK = \sqrt{OC^2 - CK^2} = \sqrt{13^2 - 12^2} = 5\text{ cm}$.\\
		Vì $K$ nằm giữa $O$ và $A$ nên $AK = OA - OK = 13 - 5 = 8\text{ cm}$.
	}
\end{bt}

% --- 9H5?14-2: Góc ở tâm, số đo cung ---

\begin{ex}%[9H5N14-2]
% nguon: SGK Toán 9 KNTT tập 1 – Bài 14, Cung và dây của một đường tròn
	Góc ở tâm chắn nửa đường tròn có số đo bằng
	\choice
	{$60^\circ$}
	{$90^\circ$}
	{\True $180^\circ$}
	{$360^\circ$}
	\loigiai{
		Nửa đường tròn có số đo bằng $180^\circ$, nên góc ở tâm chắn nửa đường tròn là góc bẹt có số đo $180^\circ$.
	}
\end{ex}

\begin{ex}%[9H5N14-2]
% nguon: SGK Toán 9 KNTT tập 1 – Bài 14, Cung và dây của một đường tròn
	Trên đường tròn $(O)$, hai điểm $M, N$ chia đường tròn thành hai cung, trong đó cung nhỏ có số đo $110^\circ$. Số đo cung lớn $\overset{\frown}{MN}$ bằng
	\choice
	{$70^\circ$}
	{$110^\circ$}
	{$220^\circ$}
	{\True $250^\circ$}
	\loigiai{
		Số đo cung lớn $\overset{\frown}{MN}$ bằng $360^\circ - 110^\circ = 250^\circ$.
	}
\end{ex}

\begin{ex}%[9H5H14-2]
% nguon: SGK Toán 9 KNTT tập 1 – Bài 14, Cung và dây của một đường tròn
	Trên đường tròn $(O; 6\text{ cm})$, lấy hai điểm $A, B$ sao cho $\widehat{AOB} = 90^\circ$. Xét tính đúng, sai của các khẳng định sau:
	\choiceTF
	{\True Số đo cung nhỏ $\overset{\frown}{AB}$ bằng $90^\circ$}
	{\True Số đo cung lớn $\overset{\frown}{AB}$ bằng $270^\circ$}
	{Tam giác $OAB$ là tam giác đều}
	{\True Độ dài dây $AB$ bằng $6\sqrt{2}\text{ cm}$}
	\loigiai{
		\begin{itemchoice}
			\itemch Đúng. Số đo cung nhỏ $\overset{\frown}{AB}$ bằng số đo góc ở tâm $\widehat{AOB} = 90^\circ$.
			\itemch Đúng. Số đo cung lớn $\overset{\frown}{AB}$ bằng $360^\circ - 90^\circ = 270^\circ$.
			\itemch Sai. Vì $OA = OB = 6\text{ cm}$ và $\widehat{AOB} = 90^\circ$ nên tam giác $OAB$ là tam giác vuông cân tại $O$.
			\itemch Đúng. Theo định lí Pythagore trong tam giác vuông $OAB$: $AB = \sqrt{6^2 + 6^2} = 6\sqrt{2}\text{ cm}$.
		\end{itemchoice}
	}
\end{ex}

\begin{ex}%[9H5H14-2]
% nguon: SGK Toán 9 KNTT tập 1 – Bài 14, Cung và dây của một đường tròn
	Trên đường tròn $(O; 5\text{ cm})$, lấy ba điểm $A, B, C$ theo thứ tự đó sao cho số đo cung nhỏ $\overset{\frown}{AB}$ bằng $60^\circ$ và số đo cung nhỏ $\overset{\frown}{BC}$ bằng $60^\circ$. Xét tính đúng, sai của các khẳng định sau:
	\choiceTF
	{\True $\widehat{AOB} = 60^\circ$}
	{\True Số đo cung nhỏ $\overset{\frown}{AC}$ bằng $120^\circ$}
	{Độ dài dây $AB$ bằng $10\text{ cm}$}
	{\True Tứ giác $OABC$ là hình thoi có diện tích bằng $\dfrac{25\sqrt{3}}{2}\text{ cm}^2$}
	\loigiai{
		\begin{itemchoice}
			\itemch Đúng. Góc ở tâm $\widehat{AOB}$ bằng số đo cung nhỏ $\overset{\frown}{AB} = 60^\circ$.
			\itemch Đúng. Vì $B$ nằm trên cung nhỏ $\overset{\frown}{AC}$ nên số đo cung nhỏ $\overset{\frown}{AC}$ bằng $60^\circ + 60^\circ = 120^\circ$.
			\itemch Sai. Tam giác $OAB$ cân tại $O$ có $\widehat{AOB} = 60^\circ$ nên là tam giác đều, suy ra $AB = OA = 5\text{ cm}$.
			\itemch Đúng. Hai tam giác $OAB$ và $OBC$ đều cạnh $5\text{ cm}$ nên $OA = AB = BC = CO = 5\text{ cm}$, do đó $OABC$ là hình thoi có diện tích $2 \cdot \dfrac{5^2\sqrt{3}}{4} = \dfrac{25\sqrt{3}}{2}\text{ cm}^2$.
		\end{itemchoice}
	}
\end{ex}

\begin{bt}%[9H5V14-2]
% nguon: SGK Toán 9 KNTT tập 1 – Bài 14, Cung và dây của một đường tròn
	Lúc $5$ giờ đúng, kim giờ và kim phút của một đồng hồ tạo thành góc ở tâm chắn cung nhỏ trên mặt đồng hồ có số đo bằng bao nhiêu độ?
	\shortans{150}
	\loigiai{
		Mỗi khoảng giờ trên mặt đồng hồ ứng với góc ở tâm $\dfrac{360^\circ}{12} = 30^\circ$.\\
		Lúc $5$ giờ đúng, kim phút chỉ số $12$ và kim giờ chỉ số $5$ nên góc ở tâm chắn cung nhỏ có số đo bằng $5 \cdot 30^\circ = 150^\circ$.
	}
\end{bt}

\begin{bt}%[9H5V14-2]
% nguon: SGK Toán 9 KNTT tập 1 – Bài 14, Cung và dây của một đường tròn
	Trên đường tròn $(O)$, hai điểm $A, B$ chia đường tròn thành hai cung sao cho số đo cung lớn gấp $3$ lần số đo cung nhỏ. Tính số đo cung nhỏ $\overset{\frown}{AB}$ (đơn vị độ).
	\shortans{90}
	\loigiai{
		Gọi $x$ (độ) là số đo cung nhỏ $\overset{\frown}{AB}$ ($0 < x < 180$). Khi đó số đo cung lớn là $3x$.\\
		Ta có $x + 3x = 360 \Rightarrow 4x = 360 \Rightarrow x = 90$. Vậy số đo cung nhỏ $\overset{\frown}{AB}$ bằng $90^\circ$.
	}
\end{bt}

% --- 9H5?14-3: Liên hệ giữa dây và khoảng cách đến tâm ---

\begin{ex}%[9H5N14-3]
% nguon: SGK Toán 9 KNTT tập 1 – Bài 14, Cung và dây của một đường tròn
	Trong một đường tròn, hai dây bằng nhau khi và chỉ khi
	\choice
	{\True Hai dây đó cách đều tâm}
	{Hai dây đó song song với nhau}
	{Hai dây đó vuông góc với nhau}
	{Hai dây đó cắt nhau tại tâm}
	\loigiai{
		Theo tính chất liên hệ giữa dây và khoảng cách từ tâm đến dây, trong một đường tròn, hai dây bằng nhau khi và chỉ khi chúng cách đều tâm.
	}
\end{ex}

\begin{ex}%[9H5N14-3]
% nguon: SGK Toán 9 KNTT tập 1 – Bài 14, Cung và dây của một đường tròn
	Cho đường tròn $(O; 10\text{ cm})$ và dây $CD = 12\text{ cm}$. Khoảng cách từ tâm $O$ đến dây $CD$ bằng
	\choice
	{$4\text{ cm}$}
	{$6\text{ cm}$}
	{\True $8\text{ cm}$}
	{$9\text{ cm}$}
	\loigiai{
		Gọi $K$ là trung điểm dây $CD$, ta có $OK \perp CD$ và $CK = 6\text{ cm}$.\\
		Trong tam giác $OKC$ vuông tại $K$: $OK = \sqrt{10^2 - 6^2} = 8\text{ cm}$.
	}
\end{ex}

\begin{ex}%[9H5H14-3]
% nguon: SGK Toán 9 KNTT tập 1 – Bài 14, Cung và dây của một đường tròn
	Cho đường tròn $(O; 13\text{ cm})$ có hai dây $AB = 10\text{ cm}$ và $CD = 24\text{ cm}$. Xét tính đúng, sai của các khẳng định sau:
	\choiceTF
	{\True Khoảng cách từ tâm $O$ đến dây $AB$ bằng $12\text{ cm}$}
	{\True Khoảng cách từ tâm $O$ đến dây $CD$ bằng $5\text{ cm}$}
	{Dây $AB$ gần tâm $O$ hơn dây $CD$}
	{\True Nếu $AB \parallel CD$ và hai dây nằm khác phía đối với tâm $O$ thì khoảng cách giữa hai dây bằng $17\text{ cm}$}
	\loigiai{
		\begin{itemchoice}
			\itemch Đúng. Khoảng cách từ $O$ đến $AB$ là $\sqrt{13^2 - 5^2} = 12\text{ cm}$.
			\itemch Đúng. Khoảng cách từ $O$ đến $CD$ là $\sqrt{13^2 - 12^2} = 5\text{ cm}$.
			\itemch Sai. Vì $5 < 12$ nên dây $CD$ gần tâm $O$ hơn dây $AB$ (dây lớn hơn thì gần tâm hơn).
			\itemch Đúng. Khi hai dây song song nằm khác phía đối với tâm $O$, khoảng cách giữa chúng bằng tổng hai khoảng cách từ tâm đến mỗi dây: $12 + 5 = 17\text{ cm}$.
		\end{itemchoice}
	}
\end{ex}

\begin{ex}%[9H5H14-3]
% nguon: SGK Toán 9 KNTT tập 1 – Bài 14, Cung và dây của một đường tròn
	Cho đường tròn $(O; 5\text{ cm})$ và điểm $M$ nằm trong đường tròn sao cho $OM = 3\text{ cm}$. Xét các dây của đường tròn đi qua điểm $M$:
	\choiceTF
	{\True Dây lớn nhất đi qua $M$ có độ dài bằng $10\text{ cm}$}
	{\True Dây ngắn nhất đi qua $M$ vuông góc với $OM$ tại $M$}
	{Độ dài dây ngắn nhất đi qua $M$ bằng $6\text{ cm}$}
	{\True Diện tích tam giác tạo bởi tâm $O$ và hai đầu mút của dây ngắn nhất đi qua $M$ bằng $12\text{ cm}^2$}
	\loigiai{
		\begin{itemchoice}
			\itemch Đúng. Dây lớn nhất đi qua $M$ là đường kính đi qua $M$, có độ dài $2R = 10\text{ cm}$.
			\itemch Đúng. Khoảng cách từ $O$ đến mọi dây đi qua $M$ đều không vượt quá $OM$, đạt lớn nhất bằng $OM$ khi dây vuông góc với $OM$ tại $M$, khi đó dây có độ dài nhỏ nhất.
			\itemch Sai. Độ dài dây ngắn nhất là $2\sqrt{5^2 - 3^2} = 2 \cdot 4 = 8\text{ cm}$.
			\itemch Đúng. Diện tích tam giác tạo bởi $O$ và dây ngắn nhất là $\dfrac{1}{2} \cdot 8 \cdot 3 = 12\text{ cm}^2$.
		\end{itemchoice}
	}
\end{ex}

\begin{bt}%[9H5V14-3]
% nguon: SGK Toán 9 KNTT tập 1 – Bài 14, Cung và dây của một đường tròn
	Cho đường tròn $(O; 17\text{ cm})$ có hai dây song song $AB = 16\text{ cm}$ và $CD = 30\text{ cm}$ nằm về hai phía đối với tâm $O$. Tính khoảng cách giữa hai dây $AB$ và $CD$ (đơn vị cm).
	\shortans{23}
	\loigiai{
		Khoảng cách từ tâm $O$ đến dây $AB$ là $d_1 = \sqrt{17^2 - 8^2} = \sqrt{225} = 15\text{ cm}$.\\
		Khoảng cách từ tâm $O$ đến dây $CD$ là $d_2 = \sqrt{17^2 - 15^2} = \sqrt{64} = 8\text{ cm}$.\\
		Vì hai dây song song nằm về hai phía của tâm $O$ nên khoảng cách giữa chúng là $d_1 + d_2 = 15 + 8 = 23\text{ cm}$.
	}
\end{bt}

\begin{bt}%[9H5V14-3]
% nguon: SGK Toán 9 KNTT tập 1 – Bài 14, Cung và dây của một đường tròn
	Một cổng vòm có dạng cung tròn bán kính $5\text{ m}$, bề rộng chân vòm (dây cung trên mặt đất) là $8\text{ m}$ và tâm đường tròn nằm dưới mặt đất. Tính chiều cao lớn nhất của cổng vòm tính từ mặt đất đến đỉnh vòm (đơn vị mét).
	\shortans{2}
	\loigiai{
		Khoảng cách từ tâm đường tròn đến dây cung mặt đất là $d = \sqrt{5^2 - 4^2} = 3\text{ m}$.\\
		Chiều cao lớn nhất từ mặt đất đến đỉnh vòm bằng $R - d = 5 - 3 = 2\text{ m}$.
	}
\end{bt}

% --- 9H5?15-1: Độ dài đường tròn, cung tròn ---

\begin{ex}%[9H5N15-1]
% nguon: SGK Toán 9 KNTT tập 1 – Bài 15, Độ dài của cung tròn. Diện tích hình quạt tròn và hình vành khuyên
	Công thức tính độ dài $l$ của cung tròn $n^\circ$ trên đường tròn bán kính $R$ là
	\choice
	{\True $l = \dfrac{n}{180}\pi R$}
	{$l = \dfrac{n}{360}\pi R$}
	{$l = \dfrac{n}{180}\pi R^2$}
	{$l = \dfrac{n}{360}\pi R^2$}
	\loigiai{
		Trên đường tròn bán kính $R$, độ dài cung tròn có số đo $n^\circ$ được tính theo công thức $l = \dfrac{n}{180}\pi R$.
	}
\end{ex}

\begin{ex}%[9H5N15-1]
% nguon: SGK Toán 9 KNTT tập 1 – Bài 15, Độ dài của cung tròn. Diện tích hình quạt tròn và hình vành khuyên
	Chu vi của đường tròn có đường kính $d = 10\text{ cm}$ bằng
	\choice
	{$5\pi\text{ cm}$}
	{\True $10\pi\text{ cm}$}
	{$20\pi\text{ cm}$}
	{$25\pi\text{ cm}$}
	\loigiai{
		Chu vi đường tròn có đường kính $d = 10\text{ cm}$ là $C = \pi d = 10\pi\text{ cm}$.
	}
\end{ex}

\begin{ex}%[9H5H15-1]
% nguon: SGK Toán 9 KNTT tập 1 – Bài 15, Độ dài của cung tròn. Diện tích hình quạt tròn và hình vành khuyên
	Cho đường tròn $(O; 18\text{ cm})$ và cung nhỏ $\overset{\frown}{MN}$ có số đo $40^\circ$. Xét tính đúng, sai của các khẳng định sau:
	\choiceTF
	{\True Độ dài đường tròn $(O)$ bằng $36\pi\text{ cm}$}
	{\True Độ dài cung nhỏ $\overset{\frown}{MN}$ bằng $4\pi\text{ cm}$}
	{Độ dài cung lớn $\overset{\frown}{MN}$ bằng $30\pi\text{ cm}$}
	{\True Nếu giữ nguyên bán kính và tăng số đo cung lên $120^\circ$ thì độ dài cung mới gấp $3$ lần độ dài cung nhỏ $\overset{\frown}{MN}$}
	\loigiai{
		\begin{itemchoice}
			\itemch Đúng. Độ dài đường tròn là $C = 2\pi \cdot 18 = 36\pi\text{ cm}$.
			\itemch Đúng. Độ dài cung nhỏ $\overset{\frown}{MN}$ là $l = \dfrac{40}{180}\pi \cdot 18 = 4\pi\text{ cm}$.
			\itemch Sai. Độ dài cung lớn $\overset{\frown}{MN}$ là $36\pi - 4\pi = 32\pi\text{ cm}$.
			\itemch Đúng. Với bán kính không đổi, độ dài cung tròn tỉ lệ thuận với số đo cung; vì $120^\circ = 3 \cdot 40^\circ$ nên độ dài cung mới gấp $3$ lần.
		\end{itemchoice}
	}
\end{ex}

\begin{ex}%[9H5H15-1]
% nguon: SGK Toán 9 KNTT tập 1 – Bài 15, Độ dài của cung tròn. Diện tích hình quạt tròn và hình vành khuyên
	Một cung tròn trên đường tròn bán kính $R = 15\text{ cm}$ có độ dài bằng $5\pi\text{ cm}$. Xét tính đúng, sai của các khẳng định sau:
	\choiceTF
	{\True Chu vi đường tròn chứa cung đó bằng $30\pi\text{ cm}$}
	{\True Số đo của cung tròn đó bằng $60^\circ$}
	{Độ dài dây căng cung tròn đó bằng $10\text{ cm}$}
	{\True Chu vi hình quạt tròn tạo bởi cung đó (lấy $\pi \approx 3{,}14$) bằng $45{,}7\text{ cm}$}
	\loigiai{
		\begin{itemchoice}
			\itemch Đúng. Chu vi đường tròn là $C = 2\pi \cdot 15 = 30\pi\text{ cm}$.
			\itemch Đúng. Từ $l = \dfrac{n}{180}\pi R \Rightarrow 5\pi = \dfrac{n}{180}\pi \cdot 15 \Rightarrow n = 60^\circ$.
			\itemch Sai. Góc ở tâm bằng $60^\circ$ nên tam giác tạo bởi tâm và hai mút cung là tam giác đều cạnh $15\text{ cm}$, suy ra dây dài $15\text{ cm}$.
			\itemch Đúng. Chu vi hình quạt tròn bằng $2R + l = 30 + 5\pi \approx 30 + 5 \cdot 3{,}14 = 45{,}7\text{ cm}$.
		\end{itemchoice}
	}
\end{ex}

\begin{bt}%[9H5V15-1]
% nguon: SGK Toán 9 KNTT tập 1 – Bài 15, Độ dài của cung tròn. Diện tích hình quạt tròn và hình vành khuyên
	Một chiếc đu quay có bán kính $20\text{ m}$. Khi đu quay quay một góc $90^\circ$ thì một cabin trên vành đu quay di chuyển được quãng đường dài bao nhiêu mét? (Lấy $\pi \approx 3{,}14$, làm tròn kết quả đến hàng phần mười).
	\shortans{31,4}
	\loigiai{
		Quãng đường cabin di chuyển là độ dài cung tròn $90^\circ$ bán kính $R = 20\text{ m}$:
		$$l = \dfrac{90}{180}\pi \cdot 20 = 10\pi \approx 10 \cdot 3{,}14 = 31{,}4\text{ m}.$$
	}
\end{bt}

\begin{bt}%[9H5V15-1]
% nguon: SGK Toán 9 KNTT tập 1 – Bài 15, Độ dài của cung tròn. Diện tích hình quạt tròn và hình vành khuyên
	Một cung tròn có số đo $72^\circ$ và độ dài bằng $6\pi\text{ cm}$. Tính bán kính của đường tròn chứa cung đó (đơn vị cm).
	\shortans{15}
	\loigiai{
		Áp dụng công thức $l = \dfrac{n}{180}\pi R$, ta có:
		$$6\pi = \dfrac{72}{180}\pi R \Rightarrow \dfrac{2}{5}R = 6 \Rightarrow R = 15\text{ cm}.$$
	}
\end{bt}

% --- 9H5?15-2: Diện tích hình tròn, hình quạt tròn ---

\begin{ex}%[9H5N15-2]
% nguon: SGK Toán 9 KNTT tập 1 – Bài 15, Độ dài của cung tròn. Diện tích hình quạt tròn và hình vành khuyên
	Công thức tính diện tích $S$ của hình quạt tròn bán kính $R$, ứng với cung có số đo $n^\circ$ là
	\choice
	{$S = \dfrac{n}{180}\pi R$}
	{$S = \dfrac{n}{180}\pi R^2$}
	{\True $S = \dfrac{n}{360}\pi R^2$}
	{$S = \dfrac{n}{360}\pi R$}
	\loigiai{
		Diện tích hình quạt tròn bán kính $R$, cung $n^\circ$ được tính theo công thức $S = \dfrac{n}{360}\pi R^2$.
	}
\end{ex}

\begin{ex}%[9H5N15-2]
% nguon: SGK Toán 9 KNTT tập 1 – Bài 15, Độ dài của cung tròn. Diện tích hình quạt tròn và hình vành khuyên
	Diện tích của hình tròn có đường kính $d = 8\text{ cm}$ bằng
	\choice
	{$8\pi\text{ cm}^2$}
	{\True $16\pi\text{ cm}^2$}
	{$32\pi\text{ cm}^2$}
	{$64\pi\text{ cm}^2$}
	\loigiai{
		Bán kính hình tròn là $R = \dfrac{8}{2} = 4\text{ cm}$, suy ra diện tích $S = \pi R^2 = 16\pi\text{ cm}^2$.
	}
\end{ex}

\begin{ex}%[9H5H15-2]
% nguon: SGK Toán 9 KNTT tập 1 – Bài 15, Độ dài của cung tròn. Diện tích hình quạt tròn và hình vành khuyên
	Cho hình quạt tròn bán kính $R = 12\text{ cm}$, ứng với cung có số đo $30^\circ$. Xét tính đúng, sai của các khẳng định sau:
	\choiceTF
	{\True Diện tích hình tròn $(O; 12\text{ cm})$ bằng $144\pi\text{ cm}^2$}
	{\True Diện tích hình quạt tròn bằng $12\pi\text{ cm}^2$}
	{Độ dài cung của hình quạt tròn bằng $4\pi\text{ cm}$}
	{\True Tỉ số giữa diện tích hình quạt tròn và diện tích hình tròn bằng $\dfrac{1}{12}$}
	\loigiai{
		\begin{itemchoice}
			\itemch Đúng. Diện tích hình tròn là $\pi \cdot 12^2 = 144\pi\text{ cm}^2$.
			\itemch Đúng. Diện tích hình quạt tròn là $\dfrac{30}{360}\pi \cdot 12^2 = 12\pi\text{ cm}^2$.
			\itemch Sai. Độ dài cung tròn là $\dfrac{30}{180}\pi \cdot 12 = 2\pi\text{ cm}$.
			\itemch Đúng. Tỉ số diện tích bằng $\dfrac{12\pi}{144\pi} = \dfrac{1}{12}$.
		\end{itemchoice}
	}
\end{ex}

\begin{ex}%[9H5H15-2]
% nguon: SGK Toán 9 KNTT tập 1 – Bài 15, Độ dài của cung tròn. Diện tích hình quạt tròn và hình vành khuyên
	Một chiếc bánh pizza hình tròn đường kính $30\text{ cm}$ được cắt đều thành $6$ miếng hình quạt tròn bằng nhau. Xét tính đúng, sai của các khẳng định sau:
	\choiceTF
	{\True Bán kính chiếc bánh bằng $15\text{ cm}$}
	{\True Góc ở tâm của mỗi miếng bánh bằng $60^\circ$}
	{Diện tích bề mặt mỗi miếng bánh bằng $45\pi\text{ cm}^2$}
	{\True Diện tích bề mặt của $2$ miếng bánh (lấy $\pi \approx 3{,}14$, làm tròn đến hàng đơn vị) bằng $236\text{ cm}^2$}
	\loigiai{
		\begin{itemchoice}
			\itemch Đúng. Bán kính chiếc bánh là $R = \dfrac{30}{2} = 15\text{ cm}$.
			\itemch Đúng. Góc ở tâm mỗi miếng bánh bằng $\dfrac{360^\circ}{6} = 60^\circ$.
			\itemch Sai. Diện tích mỗi miếng bánh là $\dfrac{1}{6}\pi \cdot 15^2 = 37{,}5\pi\text{ cm}^2$.
			\itemch Đúng. Diện tích $2$ miếng bánh là $2 \cdot 37{,}5\pi = 75\pi \approx 75 \cdot 3{,}14 = 235{,}5 \approx 236\text{ cm}^2$.
		\end{itemchoice}
	}
\end{ex}

\begin{bt}%[9H5V15-2]
% nguon: SGK Toán 9 KNTT tập 1 – Bài 15, Độ dài của cung tròn. Diện tích hình quạt tròn và hình vành khuyên
	Một hình quạt tròn có bán kính $10\text{ cm}$ và độ dài cung bằng $15\text{ cm}$. Tính diện tích của hình quạt tròn đó (đơn vị $\text{cm}^2$).
	\shortans{75}
	\loigiai{
		Áp dụng công thức liên hệ giữa diện tích hình quạt tròn và độ dài cung:
		$$S = \dfrac{l \cdot R}{2} = \dfrac{15 \cdot 10}{2} = 75\text{ cm}^2.$$
	}
\end{bt}

\begin{bt}%[9H5V15-2]
% nguon: SGK Toán 9 KNTT tập 1 – Bài 15, Độ dài của cung tròn. Diện tích hình quạt tròn và hình vành khuyên
	Một cần gạt nước kính chắn gió ô tô quét qua một góc $120^\circ$, tạo thành hình quạt tròn bán kính $30\text{ cm}$. Tính diện tích phần mặt kính được quét (đơn vị $\text{cm}^2$, lấy $\pi \approx 3{,}14$, làm tròn kết quả đến hàng đơn vị).
	\shortans{942}
	\loigiai{
		Diện tích phần kính được quét là diện tích hình quạt tròn bán kính $R = 30\text{ cm}$, cung $120^\circ$:
		$$S = \dfrac{120}{360}\pi \cdot 30^2 = 300\pi \approx 300 \cdot 3{,}14 = 942\text{ cm}^2.$$
	}
\end{bt}

% --- 9H5?15-3: Diện tích hình vành khuyên, bài toán thực tế ---

\begin{ex}%[9H5N15-3]
% nguon: SGK Toán 9 KNTT tập 1 – Bài 15, Độ dài của cung tròn. Diện tích hình quạt tròn và hình vành khuyên
	Công thức tính diện tích $S$ của hình vành khuyên giới hạn bởi hai đường tròn đồng tâm bán kính $R$ và $r$ ($R > r$) là
	\choice
	{$S = \pi(R - r)^2$}
	{\True $S = \pi(R^2 - r^2)$}
	{$S = 2\pi(R - r)$}
	{$S = \pi(R^2 + r^2)$}
	\loigiai{
		Diện tích hình vành khuyên bằng hiệu diện tích hai hình tròn đồng tâm: $S = \pi R^2 - \pi r^2 = \pi(R^2 - r^2)$.
	}
\end{ex}

\begin{ex}%[9H5N15-3]
% nguon: SGK Toán 9 KNTT tập 1 – Bài 15, Độ dài của cung tròn. Diện tích hình quạt tròn và hình vành khuyên
	Diện tích hình vành khuyên tạo bởi hai đường tròn đồng tâm có bán kính $R = 6\text{ cm}$ và $r = 4\text{ cm}$ bằng
	\choice
	{$4\pi\text{ cm}^2$}
	{$10\pi\text{ cm}^2$}
	{\True $20\pi\text{ cm}^2$}
	{$52\pi\text{ cm}^2$}
	\loigiai{
		Diện tích hình vành khuyên là $S = \pi(6^2 - 4^2) = \pi(36 - 16) = 20\pi\text{ cm}^2$.
	}
\end{ex}

\begin{ex}%[9H5H15-3]
% nguon: SGK Toán 9 KNTT tập 1 – Bài 15, Độ dài của cung tròn. Diện tích hình quạt tròn và hình vành khuyên
	Cho hai đường tròn đồng tâm có bán kính $R = 10\text{ cm}$ và $r = 6\text{ cm}$. Xét tính đúng, sai của các khẳng định sau:
	\choiceTF
	{\True Diện tích hình tròn lớn bằng $100\pi\text{ cm}^2$}
	{\True Diện tích hình tròn nhỏ bằng $36\pi\text{ cm}^2$}
	{Diện tích hình vành khuyên giới hạn bởi hai đường tròn bằng $16\pi\text{ cm}^2$}
	{\True Bán kính của hình tròn có diện tích bằng diện tích hình vành khuyên trên là $8\text{ cm}$}
	\loigiai{
		\begin{itemchoice}
			\itemch Đúng. Diện tích hình tròn lớn là $\pi \cdot 10^2 = 100\pi\text{ cm}^2$.
			\itemch Đúng. Diện tích hình tròn nhỏ là $\pi \cdot 6^2 = 36\pi\text{ cm}^2$.
			\itemch Sai. Diện tích hình vành khuyên là $100\pi - 36\pi = 64\pi\text{ cm}^2$.
			\itemch Đúng. Gọi $R_0$ là bán kính hình tròn có diện tích $64\pi\text{ cm}^2$, ta có $\pi R_0^2 = 64\pi \Rightarrow R_0 = 8\text{ cm}$.
		\end{itemchoice}
	}
\end{ex}

\begin{ex}%[9H5H15-3]
% nguon: SGK Toán 9 KNTT tập 1 – Bài 15, Độ dài của cung tròn. Diện tích hình quạt tròn và hình vành khuyên
	Một hồ nước hình tròn có đường kính $20\text{ m}$. Bao quanh hồ là một con đường đi bộ hình vành khuyên rộng $2\text{ m}$. Xét tính đúng, sai của các khẳng định sau:
	\choiceTF
	{\True Bán kính của hồ nước bằng $10\text{ m}$}
	{\True Bán kính vòng ngoài của con đường bằng $12\text{ m}$}
	{Diện tích con đường đi bộ bằng $40\pi\text{ m}^2$}
	{\True Diện tích con đường đi bộ (lấy $\pi \approx 3{,}14$, làm tròn đến hàng phần mười) bằng $138{,}2\text{ m}^2$}
	\loigiai{
		\begin{itemchoice}
			\itemch Đúng. Bán kính hồ nước là $r = \dfrac{20}{2} = 10\text{ m}$.
			\itemch Đúng. Bán kính vòng ngoài con đường là $R = 10 + 2 = 12\text{ m}$.
			\itemch Sai. Diện tích con đường là $\pi(12^2 - 10^2) = 44\pi\text{ m}^2$.
			\itemch Đúng. Ta có $44\pi \approx 44 \cdot 3{,}14 = 138{,}16 \approx 138{,}2\text{ m}^2$.
		\end{itemchoice}
	}
\end{ex}

\begin{bt}%[9H5V15-3]
% nguon: SGK Toán 9 KNTT tập 1 – Bài 15, Độ dài của cung tròn. Diện tích hình quạt tròn và hình vành khuyên
	Một tấm biển hình tròn có bán kính $10\text{ cm}$, người ta sơn viền hình vành khuyên ở mép ngoài có bề rộng $2\text{ cm}$. Tính diện tích phần viền được sơn (đơn vị $\text{cm}^2$, lấy $\pi \approx 3{,}14$, làm tròn kết quả đến hàng đơn vị).
	\shortans{113}
	\loigiai{
		Bán kính ngoài $R = 10\text{ cm}$, bán kính phần trong không sơn viền là $r = 10 - 2 = 8\text{ cm}$.\\
		Diện tích phần viền sơn là $S = \pi(10^2 - 8^2) = 36\pi \approx 36 \cdot 3{,}14 = 113{,}04 \approx 113\text{ cm}^2$.
	}
\end{bt}

\begin{bt}%[9H5V15-3]
% nguon: SGK Toán 9 KNTT tập 1 – Bài 15, Độ dài của cung tròn. Diện tích hình quạt tròn và hình vành khuyên
	Một hình vành khuyên có diện tích bằng $24\pi\text{ cm}^2$ và bán kính đường tròn nhỏ là $r = 5\text{ cm}$. Tính bán kính $R$ của đường tròn lớn (đơn vị cm).
	\shortans{7}
	\loigiai{
		Ta có $\pi(R^2 - r^2) = 24\pi \Rightarrow R^2 - 5^2 = 24 \Rightarrow R^2 = 49 \Rightarrow R = 7\text{ cm}$.
	}
\end{bt}
